\documentclass[11pt]{article}
\usepackage[margin=1in]{geometry}
\usepackage{amsmath,amssymb,amsthm,hyperref}
\newtheorem{theorem}{Theorem}
\newtheorem{lemma}{Lemma}
\newcommand{\R}{\mathbb R}
\title{Centroids of parabola-inscribed Poncelet polygons}
\author{Principia Math In Progress}
\date{21 September 2026}
\begin{document}
\maketitle
\begin{abstract}
The vertex, perimeter, and area centroids of the focus-circle Poncelet
family trace coaxial parabolas for every polygon size. Their focal
distances are, respectively, $f/N$, $f/2$, and $f/3$, where $f$ is the
outer parabola's focal distance. The proof extends the known ellipse
centroid theorem by a controlled projective limit. In particular, it
settles Conjecture 4 of Bellio--Garcia--Reznik and its pentagon special case.
\end{abstract}

\section*{Task 1: the original problem}
Normalize the outer parabola, its focus, and the circular caustic as
\[
 \mathcal P:\ y^2+4fx=0,\qquad F=(-f,0),\qquad
 \mathcal C:\ (x+f)^2+y^2=r^2,
 \qquad 0<r<f.
\]
We consider a closed Poncelet family with $N\geq3$ distinct vertices:
successive edges touch the inner circle. This is the circumscribed
family studied in \cite{BGR}, including all its convex families in
Table 2. Points at infinity are limiting configurations, rather than
finite polygons at which a Euclidean centroid is defined.
Write $C_0,C_1,C_2$ for the centroids of the vertices, the uniformly
weighted edges, and the homogeneous polygonal lamina.
For a star traversal of a nested pair the area centroid is understood
with signed area and winding multiplicity, as in \cite{ST}; the same
proof applies. The separate intersecting-conic configurations mentioned
in the caption of Table 2 of \cite{BGR} are not circumscribed polygons
with an incircle.

\begin{theorem}\label{main}
There are real constants $v_0,v_1,v_2$, depending on the fixed family,
such that the complete finite centroid loci are
\[
 C_0:\ X=v_0-\frac{N}{4f}Y^2,\qquad
 C_1:\ X=v_1-\frac{1}{2f}Y^2,\qquad
 C_2:\ X=v_2-\frac{3}{4f}Y^2.
\]
Furthermore $v_1=\tfrac32v_2+\tfrac12f$.
Thus all three loci are nondegenerate parabolas with the same axis as
$\mathcal P$, and their focal distances are $f/N,f/2,f/3$.
\end{theorem}

The input from earlier work is the following \emph{ellipse} theorem:
for nested ellipses admitting a Poncelet family, both $C_0$ and $C_2$
trace ellipses homothetic to the outer ellipse, or single points
\cite[Theorem 1]{ST}. The removable-singularity issue for the area
centroid in that proof is addressed in \cite[Proposition 2]{CC}.
Neither the input theorem nor the tangential-polygon centroid identity
below is claimed as new here. The source retains the general parabolic
assertion as Conjecture 4 in both its current preprint and journal version.

\section*{The infinite-vertex limit}
The real projective completion of $\mathcal P$ is a circle. The Poncelet
step is a fixed-point-free, orientation-preserving circle map of order
$N$ on the chosen primitive family. One can see this by mapping the
nested pair to nested ellipses. Consequently every orbit has $N$
distinct vertices, including the orbit through the unique point
$I$ of $\mathcal P$ at infinity.

Let one vertex be
\[
 Z(t)=(-t^2/(4f),t),\qquad t\longrightarrow\pm\infty.
\]
All other vertices tend to finite, distinct limits. The two edges
incident with $I$ are the horizontal tangents $y=r$ and $y=-r$ to
$\mathcal C$. Hence the two neighboring vertices tend to
\[
 A_+=(-r^2/(4f),r),\qquad A_-=(-r^2/(4f),-r).
\]
Continuity here is projective continuity of the Poncelet map; the
other orbit points are different from $I$ and therefore remain bounded
in affine coordinates. It follows immediately that
\begin{equation}\label{vertexlimit}
 (C_0)_x=-\frac{t^2}{4fN}+O(1),\qquad
 (C_0)_y=\frac{t}{N}+O(1).
\end{equation}

For the area centroid, split the polygon into the triangle formed by
$Z(t)$ and its two neighbors and the remaining bounded polygon. This
is also a valid identity for signed areas and first moments. The
triangle's area has magnitude
\[
 \frac{r}{4f}t^2+o(t^2).
\]
The remaining polygon has bounded area and first moments. The triangle
centroid is $Z(t)/3+O(1)$, and replacing its area by the total area
changes the resulting centroid by $O(1)$. Therefore
\begin{equation}\label{arealimit}
 (C_2)_x=-\frac{t^2}{12f}+O(1),\qquad
 (C_2)_y=\frac{t}{3}+O(1).
\end{equation}
In particular, neither of the two centroid loci is contained in a
line: a relation $aX+bY+c=0$, first divided by $t^2$ and then by $t$,
would force $a=b=c=0$.

\section*{A projective limit that preserves the quadratic part}
For $\varepsilon>0$ set
\[
 T_\varepsilon(x,y)=
 \left(\frac{x}{1-\varepsilon x},
       \frac{y}{1-\varepsilon x}\right).
\]
The image of the projectively completed parabola is the ellipse
\begin{equation}\label{outer}
 \Gamma_\varepsilon:\quad
 4f\varepsilon X^2+Y^2+4fX=0.
\end{equation}
Indeed the inverse substitution is
$x=X/(1+\varepsilon X)$, $y=Y/(1+\varepsilon X)$.
The map takes the closed parabolic region $y^2+4fx\leq0$, completed
by its point at infinity, to the closed ellipse disk. The circle disk
is strictly inside that region: on $\mathcal P$ the squared distance
to $F$ equals $(f+y^2/(4f))^2\geq f^2>r^2$.
Thus $T_\varepsilon(\mathcal C)$ is an ellipse strictly inside
$\Gamma_\varepsilon$. The same projective transformation takes every
polygon of our family to a Poncelet $N$-gon for this nested pair.

Fix $k\in\{0,2\}$, and let $C_{k,\varepsilon}(P)$ be the centroid
\emph{computed after transforming all the vertices} of $P$.
We do not assert projective covariance of centroids. For each finite
polygon $P$, the vertex and area formulas give
\begin{equation}\label{continuity}
 C_{k,\varepsilon}(P)\longrightarrow C_k(P).
\end{equation}
The denominator for $C_2$ is nonzero. In the convex case this is
immediate; in a nested star family orient the edges so that the
caustic is on their left, and the fan triangles from an interior
point have positive signed area.

Choose three finite polygons for which the $C_k$ values are
noncollinear, using \eqref{vertexlimit} or \eqref{arealimit}.
Their transformed centroid values remain noncollinear for all
sufficiently small $\varepsilon>0$. The ellipse theorem then excludes
the single-point alternative and gives one equation, valid for
every polygon $P$,
\begin{equation}\label{centroidellipse}
 4f\varepsilon X^2+Y^2+
 a_\varepsilon X+b_\varepsilon Y+c_\varepsilon=0,
 \qquad (X,Y)=C_{k,\varepsilon}(P).
\end{equation}
This normalization of the quadratic terms follows precisely from
homothety to \eqref{outer}.

Substitute the three chosen centroid values into
\eqref{centroidellipse}. The coefficient matrix has rows $(X,Y,1)$;
by \eqref{continuity} it tends to an invertible matrix. The right
side, $-Y^2-4f\varepsilon X^2$, also has a finite limit. Hence
$a_\varepsilon,b_\varepsilon,c_\varepsilon$ converge to finite real
$a,b,c$. For any finite polygon, passage to the limit yields
\begin{equation}\label{parabolic}
 Y^2+aX+bY+c=0,\qquad (X,Y)=C_k(P).
\end{equation}
Divide by $t^2$ in the infinite-vertex configuration. Equations
\eqref{vertexlimit} and \eqref{arealimit} give respectively
\[
 a=\frac{4f}{N}\quad(k=0),\qquad
 a=\frac{4f}{3}\quad(k=2).
\]
In particular $a>0$, so the limiting conic is a genuine parabola.
Reflection in the $x$-axis preserves the fixed Poncelet family and
commutes with each centroid. Subtracting \eqref{parabolic} for
$(X,Y)$ and $(X,-Y)$ gives $2bY=0$ throughout the locus. Since $Y$
is unbounded, $b=0$. This proves the stated equations for $C_0,C_2$.

These equations describe the entire finite parabolas, rather than
only subarcs. Indeed, the space of unlabeled polygons is the quotient
of the real outer conic by the order-$N$ Poncelet step, hence is a
circle. Only one orbit contains $I$. Removing it leaves an interval
on which both centroids are continuous. At its two ends the single
unbounded vertex has $t\to+\infty$ and $t\to-\infty$. By
\eqref{vertexlimit} and \eqref{arealimit}, $Y$ has opposite infinite
limits. The intermediate value theorem supplies every $Y\in\R$.

\section*{The perimeter centroid}
For completeness, the identity used in \cite[Remark 1]{BGR} follows
by decomposing the polygon into triangles from $F$. If its successive
vertices are $P_i$ and its edge lengths are $\ell_i$, let
$L=\sum_i\ell_i$. Each fan triangle has area $r\ell_i/2$, so
\[
 C_2=\frac{\sum_i\ell_i(F+P_i+P_{i+1})}{3L},
 \qquad
 C_1=\frac{\sum_i\ell_i(P_i+P_{i+1})}{2L}.
\]
These formulas also hold for a nested star traversal with winding
multiplicity. Thus
\[
 C_1=F+\frac32(C_2-F).
\]
Substitution of the equation for $C_2$ gives the claimed equation for
$C_1$, with $v_1=3v_2/2+f/2$. This completes the proof.

\section*{Verification and scope}
The accompanying verifier checks the projective conic identities,
the dominant-triangle asymptotics, the centroid identity on exact
asymmetric tangential polygons, and high-precision closed families
in the $N=3,4,5,6$ cases. Numerical samples are checks of the formulas;
the limit argument above proves the assertion for all $N$.
The pentagon statement (MathDB 352453) is a special case of this result
and receives no separate campaign credit. This proof does not use the
parallel/perpendicular typo in the source's different polar conjecture.

\section*{Four proposed follow-up tasks}
\begin{enumerate}
 \item \textbf{Task 2.} Express $v_0$ and $v_2$ explicitly in terms of
 $N$, $r/f$, and standard elliptic-function invariants of the porism.
 \item \textbf{Task 3.} For a circular caustic whose center is off the
 parabola's axis, determine the axes and vertices of all three centroid
 parabolas and characterize when their axes coincide.
 \item \textbf{Task 4.} Determine, for every $N$, the vertex and area
 centroid loci of the polygons polar to this family with respect to
 $\mathcal P$, addressing Conjecture 5 of \cite{BGR}.
 \item \textbf{Task 5.} Determine the leading constant and exponential
 rate in $1-r_N/f$ as $N\to\infty$, where $r_N$ is the circular
 caustic radius admitting the convex $N$-gon family.
\end{enumerate}

\begin{thebibliography}{9}
\bibitem{BGR} F. Bellio, R. Garcia, D. Reznik,
\emph{Parabola-Inscribed Poncelet Polygons Derived from the Bicentric Family},
KoG \textbf{26} (2022), 3--15,
\href{https://doi.org/10.31896/k.26.1}{doi:10.31896/k.26.1};
\href{https://arxiv.org/abs/2111.00979}{arXiv:2111.00979v3}.
\bibitem{ST} R. Schwartz, S. Tabachnikov,
\emph{Centers of mass of Poncelet polygons, 200 years after},
Math. Intelligencer \textbf{38}(2) (2016), 29--34,
\href{https://doi.org/10.1007/s00283-016-9622-9}{doi:10.1007/s00283-016-9622-9};
\href{https://arxiv.org/abs/1607.04766}{arXiv:1607.04766}.
\bibitem{CC} A. C. Chavez-Caliz,
\emph{More about areas and centers of Poncelet polygons},
Arnold Math. J. \textbf{7} (2021), 91--106,
\href{https://doi.org/10.1007/s40598-020-00154-8}{doi:10.1007/s40598-020-00154-8};
\href{https://arxiv.org/abs/2004.05404}{arXiv:2004.05404}.
\end{thebibliography}
\end{document}
